\documentclass[11pt, a4paper]{article}

\usepackage[T1]{fontenc}
\usepackage[utf8]{inputenc}
\usepackage{tgpagella}
\usepackage[spanish, es-noshorthands]{babel}
\usepackage{amsmath, amssymb, bm}
\usepackage{tabularx}
\usepackage{booktabs}
\usepackage[most]{tcolorbox}
\usepackage[margin=2.5cm]{geometry}
\usepackage{ragged2e}
\usepackage{fancyhdr}

\pagestyle{fancy}
\fancyhf{}
\setlength{\headheight}{16pt}
\lhead{\textbf{UCB Sede Tarija} | Ing. Mecatrónica}
\chead{}
\rhead{IMT-121 Circuitos Electrónicos I | Briefing Examen EC2}
\lfoot{Prof: M.Sc. Ing. Bernardo Quiroga Turdera}
\rfoot{Página \thepage}
\renewcommand{\headrulewidth}{0.4pt}

\definecolor{ucbblue}{RGB}{0, 51, 102}

\begin{document}

\begin{tcolorbox}[colback=ucbblue!5!white, colframe=ucbblue, title=\textbf{Briefing Oficial de Evaluación: Hito 2 (Unidad EC2)}]
    \centering \textbf{Estructura, Ponderación y Reglas de la Evaluación}
\end{tcolorbox}

\section*{Arquitectura de Evaluación (Hito 2)}
La calificación del Hito 2 se estructura equilibrando el razonamiento analítico individual con la capacidad demostrada de instrumentación y redacción técnica en el laboratorio.

\begin{itemize}
    \item \textbf{50\% Examen Escrito Individual.}
    \item \textbf{50\% Componente de Laboratorio} (Integrado por el trabajo en clase, la evidencia física generada y el Informe Técnico en formato IEEE de las Experiencias A+B).
\end{itemize}

\section*{Directrices del Examen Escrito Individual}
\begin{itemize}\setlength\itemsep{0.2em}
    \item \textbf{Fecha programada:} Semana 9, Sesión 1.
    \item \textbf{Duración estipulada:} 45 a 60 minutos.
    \item \textbf{Políticas de Materiales:} 
    \begin{itemize}
        \item Examen a \textbf{libro cerrado}. \textbf{NO} se permite el uso de calculadoras ni formularios.
        \item La aritmética será deliberadamente sencilla y manejable a mano.
        \item Las expresiones o fórmulas no fundamentales que pudieran requerirse serán proporcionadas en el enunciado del problema.
    \end{itemize}
    \item \textbf{Entrega simultánea:} El mismo día del examen se debe entregar el Reporte Técnico IEEE correspondiente a las Experiencias de Laboratorio A+B.
\end{itemize}

\section*{Áreas de Evaluación y Enfoque Esperado}
El examen \textbf{no} evaluará su capacidad para memorizar fórmulas, sino su criterio para aplicar el método ingenieril: \textit{Identificar el problema $\rightarrow$ Elegir el método $\rightarrow$ Fijar referencias $\rightarrow$ Resolver $\rightarrow$ Interpretar}.

\textbf{Los temas a evaluar incluyen:}
\begin{enumerate}
    \item \textbf{Superposición:} Desactivación correcta de fuentes, manejo de polaridades y suma algebraica de contribuciones.
    \item \textbf{Equivalencias de Puerto (Thévenin/Norton):} Definición de $V_{oc}$ y $R_{th}$, aislamiento de puertos y conversión de fuentes.
    \item \textbf{Máxima Transferencia y Efecto de Carga:} Relación entre $R_L$ y $R_{th}$, diferenciación entre maximizar potencia versus maximizar tensión o eficiencia.
    \item \textbf{Interpretación de Datos o Curvas:} Capacidad de analizar una tabla de datos (ej. $V_L$ vs $R_L$), proponer hipótesis técnicas válidas frente a discrepancias y distinguir entre modelo y realidad.
\end{enumerate}

\section*{Acerca de la Microdefensa Condicional}
La microdefensa oral \textbf{no es un tercer componente de evaluación} con un porcentaje asignado, ni un castigo universal. Se trata de una auditoría técnica rápida (de 2 a 4 minutos) que se activará \textbf{exclusivamente} si existen inconsistencias graves entre la evidencia de su reporte IEEE y el rendimiento demostrado en el examen escrito. Su objetivo es clarificar la autoría y nivel de comprensión real del estudiante frente a la evidencia entregada.

\end{document}
